\textbf{Case 1:} If we only consider the frequency range where HI spectral
line can be detected by the receiver of this radio telescope,
\[ z \approx \frac{f_0 - f}{f_0} \]
\hfill[1 Mark]

The lowest frequency, 1.32\,GHz, can be used to observe the highest
redshifted HI, at
\[ z = \frac{1.42 - 1.32}{1.42} = 0.0704 \]
\hfill[1 Mark]

\textbf{Alternative solution:} Above is the cosmological redshift formula in
frequency conventionally used in radio astronomy. However, if the
cosmological redshift in optical astronomy is used, student may use
\[ z \approx \frac{f_0 - f}{f} = \frac{1.42 - 1.32}{1.32} = 0.0758 \]
and the student should be awarded full mark for this part.

\textbf{Case 2:} Use the information given to calculate the redshift limit
set by the flux limit due to furthest distance that typical HI galaxy can be
detected by the telescope

for low $z$, non-relativistic redshift
\[ d = \frac{v_r}{H_0} = \frac{cz}{H_0} \]
\hfill[1 Mark]

The HI spectral-line flux density for a galaxy with HI luminosity $L$ and
line-width $\Delta f$ at distance $d$ is
\[ \frac{L}{4\pi d^2}\times\frac{1}{\Delta f}\
   (\mathrm{W\,m^{-2}\,Hz^{-1}}) \]
\hfill[2 Mark]

Setting this to the detection limit and solve correctly, and taking special
care of converting various units
\[ S = \frac{L H_0^2}{4\pi\Delta f c^2 z^2} \ge S_{\mathrm{lim}}
     = 0.5\times10^{-26} \]
\hfill[1 Mark]

Therefore,
\[ z \le \frac{H_0}{c}\sqrt{\frac{L}{4\pi\Delta f S_{\mathrm{lim}}}}, \]
\hfill[1 Mark]

Correctly substitute numerical values,
\[ z \le \frac{67.8\,(\mathrm{km^{-1}\,Mpc^{-1}})}
     {2.99\times10^{5}\,(\mathrm{km\ s^{-1}})\times
      3.08\times10^{22}\ (\mathrm{m\ Mpc^{-1}})}
   \sqrt{\frac{10^{28}\ \mathrm{W}}
     {4\pi\times10^{6}\,(\mathrm{Hz})\times0.5\times10^{-29}\
      (\mathrm{Wm^{-2}Hz^{-1}})}} \]
% sic: the numerator's units are printed "km^-1 Mpc^-1" in the source
\hfill[1 Mark]
\[ z \le 0.0929 \]
\hfill[1 Mark]

\textbf{Conclusion:} For typical HI galaxy, the highest redshift is limited
by the receiver's lowest frequency limit and we get $z_{\max} = 0.0704$ (or
0.0758 for alternative solution given above). \hfill[1 Marks]

\emph{The last part is marked only if student calculate both instrumental
factors. If students use more precise cosmological effects, they will be
marked accordingly.}
